\section{基于运行生产资料的汽轮机系统故障知识图谱构建}
\label{sec:knowledge_graph_v5}

\subsection{引言}

燃煤机组汽轮机系统由多级汽缸、再热系统、抽汽回热系统以及凝汽器冷端等多个相互关联的热力设备组成。设备之间通过蒸汽、凝结水和给水持续进行物质与能量传递，当局部设备发生通流、汽封、阀门或换热性能异常时，故障影响往往并不局限于单一测点，而会沿实际设备接口向相邻参数传播。第二章建立的健康状态模型能够在当前运行边界下给出各状态参数的正常参考值，由此获得实测状态与健康模型输出之间的偏差信息。但状态偏差本身只能说明参数相对健康状态发生了变化，难以进一步回答异常参数属于哪个设备、不同异常之间是否存在实际热力联系以及多个异常共同指向何种故障机理等问题。

电站长期设计、运行和维护过程中积累了大量与设备状态和故障有关的生产资料，其中既包括设备结构、DCS点表和系统接口等结构化信息，也包括故障分析、检修经验、技术报告和公开文献中的非结构化知识。不同资料从设备结构、参数变化、故障原因和处置方式等角度描述同一热力过程，若能够将这些分散知识转化为统一的实体和关系，就可以为后续故障诊断提供区别于纯数据相关性的结构先验。知识图谱（Knowledge Graph，KG）能够以三元组形式组织多源异构知识，并通过图结构显式表达设备、参数、异常和故障之间的语义联系\cite{Hogan2021KnowledgeGraphs}。

因此，本章面向燃煤机组汽轮机系统构建基于运行生产资料的故障知识图谱。首先分析汽轮机系统参数之间的热力关联及典型故障特性，明确图谱需要描述的知识范围；随后整理本机DCS点表、设备输入输出关系和系统边界资料，并利用公开论文、标准及运行事件资料补充故障机理知识；在此基础上采用大语言模型生成候选实体与关系，通过领域规则、证据来源及Embedding语义对齐进行校核和标准化，形成可追溯的汽轮机系统故障知识图谱，并建立领域参数节点与实际DCS测点之间的映射关系，为下一章KG-GCN故障诊断模型提供静态结构输入。

\subsection{燃煤机组汽轮机系统参数及故障特性}
\label{subsec:turbine_fault_characteristics_v5}

汽轮机主通流沿“主蒸汽--VHP--一次再热--HP--二次再热--IP--LP--凝汽器”的方向完成蒸汽膨胀和能量转换。主蒸汽压力、温度和流量构成VHP入口热力边界，VHP排汽进入一次再热系统，一次热再蒸汽形成HP入口条件，HP排汽经二次再热后进入IP，最终蒸汽经LP膨胀后排入凝汽器。除主通流外，各级汽缸还通过抽汽支路与回热系统连接，使汽轮机本体状态与高压加热器、低压加热器、除氧器及给水系统形成进一步耦合。根据本机已核对的设备IO和DCS资料，VHP一级抽汽与1号高压加热器相连，HP二级抽汽与2号高压加热器相连，IP三级、四级抽汽分别与3号、4号高压加热器相连，LP六级、七级和八级抽汽分别与相应低压加热器支路连接。对于当前资料仍存在冲突的除氧器具体抽汽级次，本章仅保留待核状态，不在未确认情况下建立确定性拓扑关系。

汽轮机热力故障通常具有多参数耦合和沿设备结构传播的特点。固体颗粒侵蚀会改变喷嘴或叶片有效流通面积，并进一步影响关键压力分布\cite{Kubiak2002TurbineDegradation}；叶片表面沉积可造成有效通流面积减小和效率下降\cite{Kubiak2002ScaleDeposit}；级间或轴端汽封缺陷会导致蒸汽泄漏，低压侧密封异常还可能伴随空气漏入和凝汽器真空下降\cite{ZhangHu2021FaultPropagationGNN}；HP--IP中间汽封泄漏与IP效率、一次热再温度及IP排汽温度等参数具有诊断关联\cite{Xu2011HPIPLeakage}；凝汽器水侧污垢会降低换热能力并影响背压和LP可利用焓降\cite{Li2016CondenserFouling,Medica2018CondenserCondition}；给水加热器内部泄漏则会破坏正常回热过程，并在抽汽侧和给水侧形成协同状态变化\cite{Heo2012FWHLeakage,Barszcz2011FeedwaterHeater}。

由上述过程可知，一方面，不同故障可能共同作用于部分压力或温度参数，仅依靠局部数值变化容易产生故障类别混淆；另一方面，同一故障在不同参数上的响应受到设备归属和热力路径约束。因此，本章知识图谱不只记录“故障与参数相关”，还需要同时表示设备之间的介质连接、参数所属设备、故障发生位置、故障引起的异常现象以及可关联的运行参数，使后续诊断模型能够在结构化领域知识约束下处理多参数故障特征。

\subsection{基于运行生产资料的知识图谱构建方法}
\label{subsec:kg_construction_v5}

\subsubsection{汽轮机热力过程的运行生产资料}
\label{subsubsec:operation_materials_v5}

运行生产资料是故障知识图谱的知识来源。对于本文研究对象，能够直接获得的本机资料以DCS测点表、设备IO关系、系统边界核对结果及状态监测建模审计资料为主，这些资料能够确定设备名称、参数点号、参数归属、主通流方向、抽汽支路以及相邻设备接口，但无法完整覆盖全部典型故障的原因、演化过程和处置经验。为补足故障知识，进一步引入同行评议论文、设备运行维护资料、标准以及公开运行事件报告，用于提取具有通用物理意义的故障机理和异常关系。不同来源资料及其主要用途如表~\ref{tab:kg_sources_v5}所示。

\begin{table}[htbp]
    \centering
    \caption{汽轮机故障知识图谱的主要知识来源}
    \label{tab:kg_sources_v5}
    \begin{tabular}{p{3.3cm}p{5.0cm}p{4.6cm}p{1.6cm}}
        \toprule
        资料类型 & 主要内容 & 知识图谱用途 & 证据级别 \\
        \midrule
        本机DCS点表与设备IO & 测点名称、点号、设备归属、入口/出口边界 & 参数实体、设备实体及本机直接拓扑 & A0 \\
        本机建模与系统审计资料 & 主通流、抽汽、再热和设备边界核对结果 & 设备接口及DCS映射校核 & A0 \\
        同行评议论文 & 汽封泄漏、通流侵蚀、沉积、冷端和回热故障机理 & 故障--异常--参数关系 & A1 \\
        标准、官方资料与公开运行事件 & 故障现象、运行影响和处置信息 & 故障关系交叉验证与补充 & A1 \\
        热力机理推断 & 缺少直接点级证据时的响应方向假设 & 半物理故障场景辅助分析，不进入主KG邻接 & B \\
        \bottomrule
    \end{tabular}
\end{table}

为保证图谱中每条关系均可追溯，本文对知识来源设置证据等级。由本机DCS、设备IO和已核实系统拓扑直接确认的关系记为A0级；由原始论文、标准或官方运行事件直接支持的通用故障机理记为A1级；仅能依据热力规律合理推断、但缺少直接点级证据的关系记为B级。A0与A1关系可作为论文主知识图谱的组成部分，B级关系仅用于第五章半物理故障场景的机理分析，不直接作为KG-GCN主邻接矩阵中的确定边。若不同本机资料之间对同一设备关系存在冲突，则标记为A0*待核关系，在冲突消除前不进入主图和图卷积计算。

该处理方式同时区分“本机事实”和“跨机组通用机理”。例如，公开文献可以直接支持“汽封缺陷导致蒸汽泄漏”这一故障机理，但若文献并未针对本文机组给出某一具体DCS温度点的确定变化方向，则不将该点级响应写成A1确定关系。由此避免将其他机组的局部经验直接套用于本文对象。

\subsubsection{节点和关系的构建与提取}
\label{subsubsec:entity_relation_v5}

运行生产资料中的设备、参数和故障信息通常采用自然语言或工程缩写描述，同一对象在不同资料中可能存在全称、简称、系统名称和DCS点号等多种表达。为保证多源知识能够在同一图结构中统一表示，首先定义知识图谱实体和关系Schema。设汽轮机系统故障知识图谱为
\begin{equation}
G_{KG}=(V,E,\mathcal R),
\end{equation}
式中：$V$为实体集合；$E$为实体间关系边集合；$\mathcal R$为关系类型集合。根据汽轮机系统运行和故障诊断需求，将实体划分为设备、参数、异常、故障和处置五类：
\begin{equation}
V=V_{dev}\cup V_{par}\cup V_{abn}\cup V_{fault}\cup V_{act} .
\end{equation}
其中，设备实体用于表示VHP、HP、IP、LP、再热系统、抽汽支路、加热器和凝汽器等物理对象；参数实体表示可由DCS直接读取的压力、温度、流量、阀位等运行量；异常实体表示蒸汽泄漏、真空下降、通流面积变化和效率降低等状态现象；故障实体表示具有明确工程含义的故障类型；处置实体表示检查、清洗或检修等维护措施。

根据上述实体类型定义设备归属、介质流向、设备接口、发生于、导致、表现为、影响、关联监测和处置为等关系。主要关系Schema如表~\ref{tab:kg_relation_types_v5}所示。

\begin{table}[htbp]
    \centering
    \caption{汽轮机故障知识图谱主要关系类型}
    \label{tab:kg_relation_types_v5}
    \begin{tabular}{lll}
        \toprule
        关系类型 & 头实体类型 & 尾实体类型 \\
        \midrule
        归属 & 参数 & 设备 \\
        介质流向/设备接口 & 设备 & 设备 \\
        发生于 & 故障 & 设备 \\
        导致/表现为 & 故障或异常 & 异常 \\
        影响/关联监测 & 故障或异常 & 参数 \\
        处置为 & 故障 & 处置 \\
        \bottomrule
    \end{tabular}
\end{table}

在候选知识抽取阶段，首先利用大语言模型对设备资料和故障文本进行语义解析，根据预定义Schema识别设备、参数、异常、故障和处置实体，并生成“头实体--关系--尾实体”候选三元组。大语言模型只承担候选知识生成和术语理解，不直接决定候选关系是否进入最终图谱。对于每条候选关系，继续依据实体类型、关系方向、本机设备兼容性以及证据来源进行规则校验。例如，“凝汽器压力”可以作为参数节点加载实时数值，而“凝汽器压力升高”属于异常状态，不能与参数节点混为同一实体类型；若某一抽汽级次与本机设备资料不一致，则即使在通用汽轮机资料中常见，也不能作为本机确定拓扑写入A0关系。

完成规则校验后，采用Embedding语义对齐方法处理简称、同义词和不同资料中的别名表达\cite{Reimers2019SentenceBERT}。对于候选实体$v_i$和标准实体$v_j$，综合语义向量相似度与字符串相似度计算实体匹配程度：
\begin{equation}
S_{ij}=\lambda S_{emb}(v_i,v_j)+(1-\lambda)S_{str}(v_i,v_j),
\end{equation}
式中：$S_{emb}$为Embedding语义相似度；$S_{str}$为字符串相似度；$\lambda$为权重系数。仅当综合相似度达到设定阈值且实体类型一致时，候选实体才归并至标准实体，否则保留为待人工核验对象。

对于已经完成实体标准化的候选关系，再综合候选抽取、领域规则、语义对齐和证据来源形成关系置信度：
\begin{equation}
C_e=\alpha C_{llm}+\beta C_{rule}+\gamma C_{emb}+\delta C_{evi},
\end{equation}
式中：$C_{llm}$表示大语言模型候选抽取置信度；$C_{rule}$表示领域规则校验结果；$C_{emb}$表示实体语义对齐置信度；$C_{evi}$表示本机资料或文本证据强度；$\alpha$、$\beta$、$\gamma$和$\delta$为相应权重系数。对于本机结构关系，最终保留与否以DCS和设备IO核对结果为最高优先级；对于故障机理关系，则必须保留能够追溯至原始论文、标准或运行事件的来源。

通过上述过程，运行生产资料中的分散知识被转换为具有统一实体类型、关系类型和证据来源的结构化三元组。例如，本机设备资料可形成“VHP--抽汽连接--一级抽汽支路”“一级抽汽支路--连接--1号高压加热器”等拓扑三元组；公开故障资料可形成“固体颗粒侵蚀--导致--有效通流面积变化”“LP轴端汽封缺陷--导致--空气漏入低压侧”“空气漏入低压侧--导致--凝汽器真空下降”等故障知识链。多个经过校核的三元组相互连接，最终形成汽轮机系统故障知识图谱。

\subsection{知识图谱构建结果}
\label{subsec:kg_result_v5}

经实体标准化、关系校核和证据审计后，当前汽轮机系统知识库包含127个节点和123条主关系候选。其中，81条关系由本机DCS、设备IO和系统拓扑直接支持，归为A0级；40条关系由公开原始论文、标准或运行资料直接支持，归为A1级；另有2条本机资料存在冲突的A0*关系暂不进入论文主图和GCN投影。除主关系外，保留7条B级热力推断关系，仅用于半物理故障场景分析。当前规则审计未发现Blocking级Schema错误。

从设备拓扑看，图谱保存“主蒸汽--VHP--一次再热--HP--二次再热--IP--LP--凝汽器”主通流关系，并按照本机已核实资料恢复各级抽汽支路：VHP一级抽汽连接1号高压加热器，HP二级抽汽连接2号高压加热器，IP三级和四级抽汽分别连接3号和4号高压加热器，LP六级至八级抽汽分别连接相应低压加热器。对于除氧器具体抽汽级次的资料冲突，仅保留待核标记，不参与当前VHP任务子图计算。

从故障语义看，图谱将设备、故障、异常和参数连接成可追溯知识链。例如，汽封缺陷可经“蒸汽泄漏”异常进一步关联汽耗或冷端真空变化；通流部件侵蚀与有效通流面积和压力分布变化相关；凝汽器水侧污垢通过换热能力下降影响背压和LP有效焓降；回热设备泄漏则通过回热性能下降影响给水侧和机组热效率。上述故障链用于解释参数之间为什么存在诊断关联，而不是直接规定第五章半物理故障样本中每一个DCS点的固定变化幅值。

系统级知识图谱包含设备、参数、异常、故障和处置等多种语义节点，而实时KG-GCN只能向具有数值来源的参数节点加载动态特征。因此，在图谱构建完成后进一步建立标准参数概念与实际DCS测点之间的映射。第五章VHP代表性算例从系统知识图谱中抽取VHP本体、一级抽汽支路及一次再热接口相关的24个状态测点作为任务参数节点，使系统级静态领域知识能够与设备级动态状态残差在同一参数空间中建立联系。

% 图3-x：汽轮机系统故障知识图谱。最终图仅从审计后的节点/边表自动生成，不允许绘图过程新增语义关系。
% 图3-x：VHP任务参数子图及DCS测点映射。该图只表示系统KG到代表性算例参数节点的投影关系。

\subsection{本章小结}

本章面向燃煤机组汽轮机系统构建了基于运行生产资料的故障知识图谱。首先结合主通流、再热、抽汽回热和冷端过程分析参数之间的结构关联以及典型热力故障的多参数传播特征；随后以本机DCS点表、设备IO和系统边界资料确定设备及参数事实关系，并利用公开论文、标准和运行资料补充故障机理知识；在此基础上定义设备、参数、异常、故障和处置五类实体以及相应关系类型，采用大语言模型生成候选三元组，并通过领域规则、证据来源和Embedding语义对齐完成校核与标准化。最终形成具有证据可追溯性的汽轮机系统故障知识图谱，并建立标准参数概念与实际DCS状态量之间的映射，为下一章利用知识图谱约束多参数状态残差传播提供静态结构基础。